% GENERATED FILE. Edit canonical lessons, then run npm run generate:books.
% canonical-source-hash: fnv1a64:b9c62c62fa3beddf
% canonical-lessons: GU-C195-khovavu, GU-C195-rakhadvu, GU-C195-batavvu, GU-C195-thekanu, GU-C195-antar

\chapter{Lost in town}
\label{ch:gu-lost-in-town}
\hlchaptermodality{195}

\begin{chapteropening}
\textbf{By the end of this chapter:} \emph{I can say that I am lost and that someone showed me the way.}

Complete the last lesson of chapter 195: I can say that I am lost and that someone showed me the way.
\end{chapteropening}

\section[\texorpdfstring{\gu{ખોવાવું} (khovāvũ)}{ખોવાવું (khovāvũ)}]{\gu{ખોવાવું} (khovāvũ) — to get lost}

\label{lesson:GU-C195-khovavu}



\begin{quote}
Before the new one: say the Gujarati for side, then the Gujarati for to cross.
\end{quote}

\subsection*{You'll want to know: \gu{ખોવાવું}}
\textbf{\gu{ખોવાવું}} — \emph{khovāvũ} — \textquotedblleft{}to get lost\textquotedblright{}.

\begin{quote}
\textbf{\gu{હું} \gu{શહેરમાં} \gu{ખોવાઈ} \gu{ગયો}.} — \emph{hũ śahermā̃ khovāī gayo} — \textquotedblleft{}I got lost in the city.\textquotedblright{}
\end{quote}

\begin{cousinweb}
\textbf{\gu{ખોવાવું}} (\emph{khovāvũ}) is built on \emph{khovũ}, \textquotedblleft{}to lose\textquotedblright{}, with the passive ending \emph{-āvũ}. CDIAL lists \emph{khovũ}, \textquotedblleft{}to lose, waste, throw away\textquotedblright{}, under Sanskrit \emph{kṣapayati}, \textquotedblleft{}destroys, throws\textquotedblright{} (CDIAL 3651).
\end{cousinweb}

\begin{grammarlens}[title={What you've built}]
One more word for everyday life at A2.
\end{grammarlens}

\subsection*{Your turn}
\begin{itemize}
\raggedright
  \item \emph{Say it:} \emph{khovāvũ}
  \item \emph{Say it:} \emph{khovāvũ}, once more
  \item \emph{Say it:} \emph{hũ śahermā̃ khovāī gayo} — I got lost in the city.
  \item \emph{Recall:} say the Gujarati for side, then the Gujarati for to cross, then say \emph{khovāvũ} again
\end{itemize}

\begin{tcolorbox}[breakable,colback=teal!4,colframe=teal!35!black,title={Before you move on}]
What does \gu{ખોવાવું} mean? (\textquotedblleft{}To get lost\textquotedblright{}.) Say it once more.
\end{tcolorbox}
\section[\texorpdfstring{\gu{રખડવું} (rakhaḍvũ)}{રખડવું (rakhaḍvũ)}]{\gu{રખડવું} (rakhaḍvũ) — to roam, to wander about}

\label{lesson:GU-C195-rakhadvu}



\begin{quote}
Before the new one: say the Gujarati for to cross, then the Gujarati for to get lost.
\end{quote}

\subsection*{You'll want to know: \gu{રખડવું}}
\textbf{\gu{રખડવું}} — \emph{rakhaḍvũ} — \textquotedblleft{}to roam, to wander about\textquotedblright{}.

\begin{quote}
\textbf{\gu{તે} \gu{બજારમાં} \gu{રખડે} \gu{છે}.} — \emph{te bajārmā̃ rakhaḍe che} — \textquotedblleft{}He wanders about the market.\textquotedblright{}
\end{quote}

\begin{cousinweb}
The origin of \textbf{\gu{રખડવું}} (\emph{rakhaḍvũ}) is not settled. CDIAL has no entry for it, and EtymDB lists no source.
\end{cousinweb}

\begin{grammarlens}[title={What you've built}]
One more word for everyday life at A2.
\end{grammarlens}

\subsection*{Your turn}
\begin{itemize}
\raggedright
  \item \emph{Say it:} \emph{rakhaḍvũ}
  \item \emph{Say it:} \emph{rakhaḍvũ}, once more
  \item \emph{Say it:} \emph{te bajārmā̃ rakhaḍe che} — He wanders about the market.
  \item \emph{Recall:} say the Gujarati for to cross, then the Gujarati for to get lost, then say \emph{rakhaḍvũ} again
\end{itemize}

\begin{tcolorbox}[breakable,colback=teal!4,colframe=teal!35!black,title={Before you move on}]
What does \gu{રખડવું} mean? (\textquotedblleft{}To roam, to wander about\textquotedblright{}.) Say it once more.
\end{tcolorbox}
\section[\texorpdfstring{\gu{બતાવવું} (batāvvũ)}{બતાવવું (batāvvũ)}]{\gu{બતાવવું} (batāvvũ) — to show}

\label{lesson:GU-C195-batavvu}



\begin{quote}
Before the new one: say the Gujarati for to get lost, then the Gujarati for to roam.
\end{quote}

\subsection*{You'll want to know: \gu{બતાવવું}}
\textbf{\gu{બતાવવું}} — \emph{batāvvũ} — \textquotedblleft{}to show\textquotedblright{}.

\begin{quote}
\textbf{\gu{પોલીસ} \gu{મને} \gu{રસ્તો} \gu{બતાવે} \gu{છે}.} — \emph{polīs mane rasto batāve che} — \textquotedblleft{}A policeman shows me the way.\textquotedblright{}
\end{quote}

\begin{cousinweb}
The origin of \textbf{\gu{બતાવવું}} (\emph{batāvvũ}) is not settled. CDIAL lists no Gujarati form, and EtymDB lists no source.
\end{cousinweb}

\begin{grammarlens}[title={What you've built}]
One more word for everyday life at A2.
\end{grammarlens}

\subsection*{Your turn}
\begin{itemize}
\raggedright
  \item \emph{Say it:} \emph{batāvvũ}
  \item \emph{Say it:} \emph{batāvvũ}, once more
  \item \emph{Say it:} \emph{polīs mane rasto batāve che} — A policeman shows me the way.
  \item \emph{Recall:} say the Gujarati for to get lost, then the Gujarati for to roam, then say \emph{batāvvũ} again
\end{itemize}

\begin{tcolorbox}[breakable,colback=teal!4,colframe=teal!35!black,title={Before you move on}]
What does \gu{બતાવવું} mean? (\textquotedblleft{}To show\textquotedblright{}.) Say it once more.
\end{tcolorbox}
\section[\texorpdfstring{\gu{ઠેકાણું} (ṭhekāṇũ)}{ઠેકાણું (ṭhekāṇũ)}]{\gu{ઠેકાણું} (ṭhekāṇũ) — whereabouts, place}

\label{lesson:GU-C195-thekanu}



\begin{quote}
Before the new one: say the Gujarati for to roam, then the Gujarati for to show.
\end{quote}

\subsection*{You'll want to know: \gu{ઠેકાણું}}
\textbf{\gu{ઠેકાણું}} — \emph{ṭhekāṇũ} — \textquotedblleft{}whereabouts, place\textquotedblright{}.

\begin{quote}
\textbf{\gu{આ} \gu{દુકાનનું} \gu{ઠેકાણું} \gu{શું} \gu{છે}?} — \emph{ā dukānnũ ṭhekāṇũ śũ che?} — \textquotedblleft{}Where exactly is this shop?\textquotedblright{}
\end{quote}

\begin{cousinweb}
\textbf{\gu{ઠેકાણું}} (\emph{ṭhekāṇũ}) is a tadbhava. CDIAL lists it, glossed \textquotedblleft{}station\textquotedblright{}, under \emph{*ṭhēkk}, \textquotedblleft{}fix\textquotedblright{}, a reconstructed form (CDIAL 5509).
\end{cousinweb}

\begin{grammarlens}[title={What you've built}]
One more word for everyday life at A2.
\end{grammarlens}

\subsection*{Your turn}
\begin{itemize}
\raggedright
  \item \emph{Say it:} \emph{ṭhekāṇũ}
  \item \emph{Say it:} \emph{ṭhekāṇũ}, once more
  \item \emph{Say it:} \emph{ā dukānnũ ṭhekāṇũ śũ che?} — Where exactly is this shop?
  \item \emph{Recall:} say the Gujarati for to roam, then the Gujarati for to show, then say \emph{ṭhekāṇũ} again
\end{itemize}

\begin{tcolorbox}[breakable,colback=teal!4,colframe=teal!35!black,title={Before you move on}]
What does \gu{ઠેકાણું} mean? (\textquotedblleft{}Whereabouts, place\textquotedblright{}.) Say it once more.
\end{tcolorbox}
\section[\texorpdfstring{\gu{અંતર} (antar)}{અંતર (antar)}]{\gu{અંતર} (antar) — distance}

\label{lesson:GU-C195-antar}



\begin{quote}
Before the new one: say the Gujarati for to show, then the Gujarati for whereabouts.
\end{quote}

\subsection*{You'll want to know: \gu{અંતર}}
\textbf{\gu{અંતર}} — \emph{antar} — \textquotedblleft{}distance\textquotedblright{}.

\begin{quote}
\textbf{\gu{બંને} \gu{ગામ} \gu{વચ્ચે} \gu{ઘણું} \gu{અંતર} \gu{છે}.} — \emph{banne gām vacce ghaṇũ antar che} — \textquotedblleft{}The two villages are far apart.\textquotedblright{}
\end{quote}

\begin{cousinweb}
\textbf{\gu{અંતર}} (\emph{antar}) is Sanskrit \emph{ántara}, \textquotedblleft{}interior; interval\textquotedblright{}, taken over whole, a tatsama. CDIAL lists its inherited cousin \emph{ā̃tro}, \textquotedblleft{}distance, interval\textquotedblright{} (CDIAL 357).
\end{cousinweb}

\begin{grammarlens}[title={What you've built}]
One more word for everyday life at A2.
\end{grammarlens}

\subsection*{Your turn}
\begin{itemize}
\raggedright
  \item \emph{Say it:} \emph{antar}
  \item \emph{Say it:} \emph{antar}, once more
  \item \emph{Say it:} \emph{banne gām vacce ghaṇũ antar che} — The two villages are far apart.
  \item \emph{Recall:} say the Gujarati for to show, then the Gujarati for whereabouts, then say \emph{antar} again
\end{itemize}

\begin{tcolorbox}[breakable,colback=teal!4,colframe=teal!35!black,title={Before you move on}]
What does \gu{અંતર} mean? (\textquotedblleft{}Distance\textquotedblright{}.) Say it once more.
\end{tcolorbox}
